Prostorová distribuce úhrnné plodnosti v~EU27 odhaluje geografický vzorec, který není náhodný: státy s~nejvyšší plodností (FR, SE, NO, IS) jsou zároveň státy s~nejvyššími výdaji na aktivní politiku zaměstnanosti, nejsilnějším pokrytím kolektivními smlouvami a~nejrozvinutějšími institucemi slaďování práce a~rodiny.
Naopak státy jihovýchodní a~střední Evropy včetně ČR kombinují nízkou plodnost s~nízkými \ac{APZ} výdaji, slabým pokrytím \ac{KS} a~omezenou nabídkou firemních benefitů na rodičovství --- naznačujíce systémové propojení těchto fenoménů.
Pro ČR je typická kombinace silné politické rétoriky o~rodinné politice s~fakticky nízkými firemními supplementy v~\ac{KS}: zákonná mateřská a~rodičovská je relativně štědrá, ale firemní příspěvky a~záruky nad zákonné minimum jsou výjimkou (srov.~obr.~\ref{fig:natality_tfr_timeline}).
Prosazení minimálních standardů rodinných benefitů v~sektorových kolektivních smlouvách --- garantujících supplementy k~rodičovské dovolené, přednostní zkrácený úvazek nebo příspěvky na péči o~dítě --- by propojilo mzdový sociální dialog s~demografickou politikou.

Prostorová distribuce úhrnné plodnosti v~EU27 odhaluje geografický vzorec, který není náhodný: státy s~nejvyšší plodností (FR, SE, NO, IS) jsou zároveň státy s~nejvyššími výdaji na aktivní politiku zaměstnanosti, nejsilnějším pokrytím kolektivními smlouvami a~nejrozvinutějšími institucemi slaďování práce a~rodiny.
Naopak státy jihovýchodní a~střední Evropy včetně ČR kombinují nízkou plodnost s~nízkými \ac{APZ} výdaji, slabým pokrytím \ac{KS} a~omezenou nabídkou firemních benefitů na rodičovství --- naznačujíce systémové propojení těchto fenoménů.
Pro ČR je typická kombinace silné politické rétoriky o~rodinné politice s~fakticky nízkými firemními supplementy v~\ac{KS}: zákonná mateřská a~rodičovská je relativně štědrá, ale firemní příspěvky a~záruky nad zákonné minimum jsou výjimkou (srov.~obr.~\ref{fig:natality_tfr_timeline}).
Prosazení minimálních standardů rodinných benefitů v~sektorových kolektivních smlouvách --- garantujících supplementy k~rodičovské dovolené, přednostní zkrácený úvazek nebo příspěvky na péči o~dítě --- by propojilo mzdový sociální dialog s~demografickou politikou.

Prostorová distribuce úhrnné plodnosti v~EU27 odhaluje geografický vzorec, který není náhodný: státy s~nejvyšší plodností (FR, SE, NO, IS) jsou zároveň státy s~nejvyššími výdaji na aktivní politiku zaměstnanosti, nejsilnějším pokrytím kolektivními smlouvami a~nejrozvinutějšími institucemi slaďování práce a~rodiny.
Naopak státy jihovýchodní a~střední Evropy včetně ČR kombinují nízkou plodnost s~nízkými \ac{APZ} výdaji, slabým pokrytím \ac{KS} a~omezenou nabídkou firemních benefitů na rodičovství --- naznačujíce systémové propojení těchto fenoménů.
Pro ČR je typická kombinace silné politické rétoriky o~rodinné politice s~fakticky nízkými firemními supplementy v~\ac{KS}: zákonná mateřská a~rodičovská je relativně štědrá, ale firemní příspěvky a~záruky nad zákonné minimum jsou výjimkou (srov.~obr.~\ref{fig:natality_tfr_timeline}).
Prosazení minimálních standardů rodinných benefitů v~sektorových kolektivních smlouvách --- garantujících supplementy k~rodičovské dovolené, přednostní zkrácený úvazek nebo příspěvky na péči o~dítě --- by propojilo mzdový sociální dialog s~demografickou politikou.

Prostorová distribuce úhrnné plodnosti v~EU27 odhaluje geografický vzorec, který není náhodný: státy s~nejvyšší plodností (FR, SE, NO, IS) jsou zároveň státy s~nejvyššími výdaji na aktivní politiku zaměstnanosti, nejsilnějším pokrytím kolektivními smlouvami a~nejrozvinutějšími institucemi slaďování práce a~rodiny.
Naopak státy jihovýchodní a~střední Evropy včetně ČR kombinují nízkou plodnost s~nízkými \ac{APZ} výdaji, slabým pokrytím \ac{KS} a~omezenou nabídkou firemních benefitů na rodičovství --- naznačujíce systémové propojení těchto fenoménů.
Pro ČR je typická kombinace silné politické rétoriky o~rodinné politice s~fakticky nízkými firemními supplementy v~\ac{KS}: zákonná mateřská a~rodičovská je relativně štědrá, ale firemní příspěvky a~záruky nad zákonné minimum jsou výjimkou (srov.~obr.~\ref{fig:natality_tfr_timeline}).
Prosazení minimálních standardů rodinných benefitů v~sektorových kolektivních smlouvách --- garantujících supplementy k~rodičovské dovolené, přednostní zkrácený úvazek nebo příspěvky na péči o~dítě --- by propojilo mzdový sociální dialog s~demografickou politikou.

\input{texparts/python/natality_tfr_map}



